\problemstart{AP-01-010}{Content Fingerprints for Split Leakage}{Difficulty 4/5}{Chapter 1 \textbullet\ numpy}{Detect byte-identical examples crossing a train/validation boundary without relying on object identity.}

        \subsection{Mathematical statement}


Let $h$ be a collision-resistant digest over a row's dtype, shape, and canonical
C-order bytes.  Define
\[
  H(X)_i=h(\operatorname{dtype}(X_i)\,\|\,\operatorname{shape}(X_i)\,\|\,\operatorname{bytes}(X_i))
\]
and the cross-split leakage relation
\[
  \mathcal D=\{(i,j):H(X^{\mathrm{train}})_i=H(X^{\mathrm{valid}})_j\}.
\]
The diagnostic rate is $\rho=|\{j:\exists i,(i,j)\in\mathcal D\}|/N_{\mathrm{valid}}$.


        \subsection{Code problem}


Implement row fingerprints with SHA-256 from the Python standard library, then
return all sorted train/validation duplicate index pairs.  Canonicalize each row
to C-order and include dtype plus shape in the digest domain.  The function is a
data-quality diagnostic, not a cryptographic authentication protocol.


        \begin{contractbox}
        \begin{Code}
        def fingerprint_rows(array: np.ndarray) -> np.ndarray:

def cross_split_duplicates(train: np.ndarray, valid: np.ndarray) -> list[tuple[int, int]]:
        \end{Code}
        \end{contractbox}

        \subsection{Theory--engineering gap inventory}

        \begin{gapbox}
        \textbf{Explicit gap.} Disjoint index sets do not imply disjoint observed content; equality of values is separate from identity of array objects or memory addresses.

        \textbf{Implicit gap exposed by the result.} A copied row has no shared memory and a different index yet still triggers the digest relation, exposing leakage hidden by index-only checks.
        \end{gapbox}

        \subsection{Acceptance evidence}

        \begin{evidencebox}
        \begin{itemize}
          \item Known copied rows produce the exact sorted pair list, including repeated duplicates.
  \item Changing one pixel changes the digest; changing dtype also changes the digest domain.
  \item The benchmark prints duplicate-pair count, affected-validation count, and leakage rate separately. (Trust-based: the test suite does not machine-check benchmark output.)
        \end{itemize}
        \end{evidencebox}

        \subsection{Constraints}

        \begin{itemize}
          \item Use \code{hashlib.sha256}; do not invent a hash function.
  \item A Python row loop is acceptable here because hashing is a per-record boundary, not a numerical kernel.
  \item Reject rank-zero arrays and incompatible trailing row shapes with \code{ValueError}.
        \end{itemize}

        \begin{sourcebox}
        This is an atomic adaptation, not copied solution code. Nearest primary sources:
        \begin{itemize}
          \item \href{https://numpy.org/doc/stable/user/basics.copies.html}{NumPy: Copies and views}
  \item \href{https://arxiv.org/abs/1908.01242}{Kannada-MNIST paper}
        \end{itemize}
        \end{sourcebox}
